%% kapa: exam=1 round=23 year=2012 period=2 form=A subject=accounting subject_ko=회계학 seq=38 number=38
%% source: Q-Net 감정평가사 기출문제 (cst00309) artlSeq=5052580 file=raw/1008/5052580/2012년 제23회 감정평가사 제1차시험 기출문제(2교시 A형).pdf
%% conversion: kordoc 4.24.51 -> tools/gen1.py (auto); review=proofread-sonnet; flags=none
\begin{question}{38}
(주)서울은 완제품 생산에 필요한 부품을 자가제조하고 있다. 부품 10,000단위를 제조하는데 소요되는 연간제조원가는 다음과 같다.
\begin{center}
\small
\begin{tabular}{|l r|}
\hline
직접재료원가 & ￦600,000 \\
직접노무원가 & 150,000 \\
변동제조간접원가 & 50,000 \\
부품생산용설비 감가상각비 & 120,000 \\
고정제조간접원가배부액 & 70,000 \\
\cline{2-2}
\hspace{2em}총 계 & ￦990,000 \\
\hline
\end{tabular}
\end{center}
(주)서울은 (주)공덕으로부터 단위당 ￦85에 10,000단위의 부품을 공급하겠다는 제의를 받았다. 이 제의를 수락하더라도 부품생산용설비 감가상각비와 고정제조간접원가는 계속 발생한다. (주)서울이 이 제의를 수락할 경우에는 연간 ￦70,000에 설비를 임대할 수 있다. (주)서울이 이 제의를 수락하는 경우 (주)서울의 이익에 미치는 영향은? \par
\begin{choices}
  \item ￦10,000 증가
  \item ￦20,000 증가
  \item ￦20,000 감소
  \item ￦50,000 감소
  \item ￦50,000 증가
\end{choices}
\end{question}
\answer{2}
